% =====================================================================
%  Energy and Performance Analysis of ECRL and 2LAL Adiabatic Logic
%  Under Frequency, Transition-Time and Circuit-Parameter Variations
%
%  Build:   latexmk -pdf main.tex        (or pdflatex -> bibtex -> pdflatex x2)
%  Layout:  main.tex, refs.bib, figs/  (gnuplot/ngspice outputs as PDF)
%
%  Draft aids (turn off for submission: set \draftfalse below):
%    \tbd{...}    red   -> a value/result that must come from simulation
%    \guide{...}  blue  -> instruction carried over from the MS template
%    \simfig{w}{file}   -> includes file if it exists, else a labelled box,
%                          so the paper always compiles while sweeps run
%
%  DATA SOURCES (github.com/Sairam-Kesanapalli/Adiabatic_sweep):
%    ECRL : ecrl_exp1/ecrl_slowramp_sweep.csv            (as committed)
%    2LAL : 2lal_inverter_sweep_terminated.csv           (deck with the
%           chain-termination fix, 2lal_termination_fix.patch). The committed
%           2lal_inverter_sweep.csv is superseded: its stage-4 energy is
%           corrupted by the unterminated chain end (see Sec. III-D).
%    CMOS : no deck in the repo -> every CMOS number is \tbd or labelled
%           as the analytical 1/2 C V^2 reference.
% =====================================================================
\documentclass[conference]{IEEEtran}
\IEEEoverridecommandlockouts

\usepackage{cite}
\usepackage{amsmath,amssymb,amsfonts}
\usepackage{graphicx}
\usepackage{xcolor}
\usepackage{booktabs}
\usepackage{multirow}
\usepackage{array}
\newcolumntype{L}[1]{>{\raggedright\arraybackslash}p{#1}}
\usepackage[per-mode=symbol]{siunitx}
\usepackage[american]{circuitikz}
\usepackage{listings}
\usepackage{url}
\usepackage[hidelinks]{hyperref}
\usepackage[capitalise,noabbrev]{cleveref}

% ---------------- draft switches ----------------
\newif\ifdraft
\drafttrue            % <- \draftfalse hides all \guide notes and red TBD styling
\ifdraft
  \newcommand{\tbd}[1][]{\textcolor{red}{\textbf{[TBD#1]}}}
  \newcommand{\guide}[1]{\textcolor{blue}{\textit{[#1]}}}
\else
  \newcommand{\tbd}[1][]{[TBD#1]}
  \newcommand{\guide}[1]{}
\fi

% Include a figure if generated, else draw a placeholder box with the expected path.
\newlength{\simfigw}
\newcommand{\simfig}[2]{%
  \setlength{\simfigw}{#1}%
  \IfFileExists{#2}%
    {\includegraphics[width=\simfigw]{#2}}%
    {\fbox{\parbox[c][0.45\simfigw][c]{0.94\simfigw}{\centering\footnotesize
      \texttt{\detokenize{#2}}\\[2pt](not generated yet)}}}}

% ---------------- notation ----------------
\newcommand{\VDD}{V_{\mathrm{DD}}}
\newcommand{\CL}{C_{\mathrm{L}}}
\newcommand{\Tph}{T_{\mathrm{ph}}}
\newcommand{\Vth}{V_{\mathrm{th}}}
\newcommand{\Eop}{E_{\mathrm{op}}}
\newcommand{\EDP}{\mathrm{EDP}}
\newcommand{\inb}{\overline{\mathit{in}}}
\newcommand{\outb}{\overline{\mathit{out}}}

\sisetup{detect-all}
\lstset{basicstyle=\ttfamily\scriptsize, columns=fullflexible,
        frame=single, breaklines=true, keepspaces=true}

\begin{document}

\title{Energy and Performance Analysis of ECRL and 2LAL Adiabatic Logic Under
Frequency, Transition-Time and Circuit-Parameter Variations}

\author{%
\IEEEauthorblockN{Sairam Kesanapalli, Durgesh, Muthusankar Eswaran}
\IEEEauthorblockA{Department of Electronics and Communication Engineering\\
Indian Institute of Information Technology Dharwad, Karnataka, India\\
\tbd[: email addresses]}
}

\maketitle

% =====================================================================
\begin{abstract}
Conventional CMOS logic dissipates energy on every charge and discharge of its
capacitive load. Adiabatic logic reduces this loss by charging nodes from a
ramped power clock and returning the stored charge to it. This work compares
Efficient Charge Recovery Logic (ECRL) and Two-Level Adiabatic Logic (2LAL)
inverters in the SkyWater SKY130 process under one simulation and measurement
framework: ngspice, $V_{\mathrm{DD}}=1.8$~V, $L=0.15$~\textmu m, a 25~fF load on
every metered node, alternating data, and energy defined as the integral of
power-clock power over one clock period. Transition time is swept from 1 to
100~ns (10--1000~MHz transition frequency) and device width from 0.84 to
3.0~\textmu m. Across the sweep, ECRL energy per operation ranges from
11.55 to 59.59~fJ and 2LAL from 0.55 to 50.53~fJ. At $W=1$~\textmu m, 2LAL
dissipates 1.19$\times$ less than ECRL at 1~ns and 14.6$\times$ less at
100~ns. 2LAL energy falls approximately as $1/T$ (log--log slope $-0.89$ to
$-0.99$ between 10 and 50~ns), whereas ECRL flattens towards a
threshold-related residue. Channel-length (0.15--2~\textmu m) and load
(5--200~fF) sweeps confirm the adiabatic model: 2LAL energy follows
$R(\CL+C_p)^2$ with an $R$ within 8\% of the measured gate resistance, and at
a published 2LAL operating point the SKY130 ring is within 1.17$\times$ of the
reported per-FET energy. Two measurement defects that change these
conclusions, a data-static metering window and an unterminated 2LAL chain,
are identified and corrected.
\guide{Add the CMOS reduction figure once a CMOS deck is simulated; the analytical
$\tfrac12 C V^2$ reference is 40.5~fJ per transition.}
\end{abstract}

\begin{IEEEkeywords}
Adiabatic logic, ECRL, 2LAL, low-power VLSI, energy recovery, SKY130,
ngspice, energy-delay product.
\end{IEEEkeywords}

% =====================================================================
\section{Introduction}
\label{sec:intro}

\subsection{Background}
\guide{Cover: growth of low-power VLSI; dynamic switching energy; CMOS
charging/discharging; need for energy-efficient logic. Cite 5--7 references
(more or fewer as needed).}
Dynamic switching energy remains a dominant component of power in digital CMOS
\cite{chandrakasan1992,rabaey2003}. \tbd[: expand].

\subsection{Adiabatic Logic}
\guide{Concept of adiabatic switching; gradual charging; energy recovery;
effect of transition time. Include equations and define every variable.}
Adiabatic switching charges a load through a resistive path with a slowly
ramping supply, so that the voltage across the switching device stays small
throughout the transition \cite{koller1992,athas1994}. The dissipated energy
then scales with the ratio of the circuit time constant to the ramp time
(\cref{sec:adiabatic}), and the charge stored on the load can be returned to
the supply during the falling ramp.

\subsection{ECRL}
\label{sec:ecrl}
ECRL \cite{moon1996} uses four transistors and no DC supply: a cross-coupled
PMOS pair whose sources sit on the power clock $\phi$, and two NMOS pull-downs
driven by the differential input (\cref{fig:ecrl}). With $\mathit{in}$ high,
N2 holds $\mathit{out}$ at ground, P1 (gate on $\mathit{out}$) conducts, and
$\outb$ follows $\phi$ up and back down; the charge on $\outb$ returns to the
clock through P1 on the falling ramp. Thus $\mathit{out}=\overline{\mathit{in}}$
and $\outb=\mathit{in}$, and exactly one output pulses every cycle whatever the
data. The clock is a trapezoid with four equal intervals (wait, evaluate, hold,
recover); the input may change only during wait. A single gate uses one
clock phase; a cascade places successive stages on successive phases.

Its advantages are the small cell (4 FETs) and single-rail clocking per gate.
Its main limitation is incomplete recovery: on the falling ramp P1 stops
conducting once $\outb$ is within about $|V_{\mathrm{th},p}|$ of its gate, so
$\outb$ is stranded at a residual voltage (0.53~V at $\Tph=10$~ns and
0.42~V at 100~ns in our simulations). When the data changes, the NMOS dumps
that charge to ground, which is not recoverable (\cref{eq:ecrl}).

\begin{figure}[!t]
\centering
\begin{circuitikz}[scale=0.82, transform shape]
  \draw (0,4) node[pmos, xscale=-1](P1){};
  \draw (4,4) node[pmos](P2){};
  \draw (P1.S) -- (P1.S |- 0,5.2) -- (P2.S |- 0,5.2)
        node[midway, above]{$\phi$ (power clock)} -- (P2.S);
  \draw (P1.G) -- ++(0.4,0) coordinate(g1);
  \draw (P2.G) -- ++(-0.4,0) coordinate(g2);
  \draw (P1.D) -- ++(0,-1.4) coordinate(A);
  \draw (P2.D) -- ++(0,-1.4) coordinate(B);
  \draw (g1) |- (P2.D |- 0,2.85) node[circ]{};
  \draw (g2) |- (P1.D |- 0,2.45) node[circ]{};
  \draw (A) node[circ]{} -- ++(-1,0) node[left]{$\outb$};
  \draw (B) node[circ]{} -- ++(1,0) node[right]{$\mathit{out}$};
  \draw (A) -- ++(0,-0.2) node[nmos, anchor=D, xscale=-1](N1){};
  \draw (B) -- ++(0,-0.2) node[nmos, anchor=D](N2){};
  \draw (N1.G) -- ++(0.3,0) node[right]{$\inb$};
  \draw (N2.G) -- ++(-0.3,0) node[left]{$\mathit{in}$};
  \draw (N1.S) node[ground]{};
  \draw (N2.S) node[ground]{};
  \node[font=\scriptsize] at ($(P1.center)+(-0.9,0)$) {P1};
  \node[font=\scriptsize] at ($(P2.center)+(0.9,0)$) {P2};
  \node[font=\scriptsize] at ($(N1.center)+(-0.9,0)$) {N1};
  \node[font=\scriptsize] at ($(N2.center)+(0.9,0)$) {N2};
\end{circuitikz}
\caption{ECRL inverter as simulated (device names follow the netlist:
P1\,=\,XMP1, N1\,=\,XMN1). $\mathit{out}=\overline{\mathit{in}}$; a 25~fF
load on each output is omitted from the drawing.}
\label{fig:ecrl}
\end{figure}

\subsection{2LAL}
\label{sec:2lal}
2LAL \cite{frank2003mistakes,frank2020icrc} is a fully adiabatic family built
from CMOS transmission gates (TGs) and clocked by four trapezoidal rails
$\phi_0$--$\phi_3$, each one phase later than the last; stage $n$ is clocked
one phase after stage $n-1$. Data is dual-rail: a pair $(T,C)$ rests at
FALSE $=(0,V_{\mathrm{DD}})$, and TRUE is a pulse on $T$ (from $\phi_t$) while
$C$ is pulled down by the antiphase rail $\phi_{t+2}$. A logic 0 is the
\emph{absence} of a pulse; encoding it as a pulse on $C$ makes the pass-gate
pFET conduct while $\phi_t$ is high and produces spurious ones downstream.

The working cell has eight transistors (\cref{fig:2lal}): a forward pair of
TGs that drives the output pair, and a reverse pair that restores the previous
stage's nodes (from $\phi_{t-1}$ and $\phi_{t+1}$) once this stage's output is
valid. A forward-only four-transistor stage cannot pull its own output back
down and latches high. Inversion cannot be done by crossing the two rails,
because the lopsided rest state $(0,V_{\mathrm{DD}})$ becomes $(V_{\mathrm{DD}},0)$,
a permanent TRUE. It instead uses quad-rail data: two dual-rail sub-chains
carrying $A$ and $\overline{A}$, exchanged at every stage. This costs no
transistors per cell but doubles the datapath, so one inverter column is
16~FETs. Because one sub-chain always carries a TRUE, a column pulses every
period, like an ECRL gate.

Advantages: every switch obeys the adiabatic rules, so there is no threshold
residue and dissipation keeps falling with ramp time. Limitations: four clock
rails, 4$\times$ the devices of ECRL per inverter, and data-dependent floating
nodes whose correctness depends on the neighbouring stages (\cref{sec:termination}).
\guide{The cell is reconstructed from Frank's slides, a DeBenedictis report and
the 4lc PHAS2 cell; verify against the primary source if obtainable.}

\begin{figure}[!t]
\centering
\simfig{0.9\columnwidth}{figs/fig_2lal_cell.png}
\caption{2LAL stage (4 TGs, 8 FETs): forward pair drives the output pair,
reverse pair restores the previous stage. Copy
\texttt{slides\_figs/fig\_2lal\_cell.png} to \texttt{figs/}.}
\label{fig:2lal}
\end{figure}

\subsection{Research Gap}
ECRL was introduced with an energy-versus-frequency evaluation in a 1~\textmu m,
5~V process \cite{moon1996}, and 2LAL has been developed as a fully adiabatic
family by Frank and co-workers \cite{frank2003mistakes,frank2020icrc}. Those
results come from different processes, supply voltages, loads and energy
definitions, so they cannot be compared number-for-number, and an energy value
is meaningful only when the operation it covers is stated.
This work implements both families in one open PDK (SKY130) and one simulator,
and compares them under seven matched conditions: technology and corner,
supply, load, device geometry, time base, energy definition and switching
activity. It also reports two measurement defects found along the way,
because each one changed the comparison.
\guide{Add 2--3 citations of recent SKY130 or cross-family adiabatic studies
here; do not claim that no such comparison exists.}

\subsection{Objectives}
\begin{enumerate}
  \item Implement and functionally validate CMOS, ECRL and 2LAL inverters in SKY130.
  \item Quantify the effect of frequency and power-clock transition time on energy.
  \item Quantify the effect of transistor width (load capacitance: \tbd, not yet swept).
  \item Compare energy, average power, delay and energy-delay product (EDP).
  \item Identify operating regions in which each adiabatic family reduces energy
        relative to CMOS.
\end{enumerate}

% =====================================================================
\section{Theoretical Background}
\label{sec:theory}

\subsection{CMOS Switching Energy}
Charging a load $\CL$ from $0$ to $\VDD$ through a constant supply draws charge
$\CL\VDD$ and therefore energy $\CL\VDD^2$ from the supply. Half of it is stored
on the load; the other half is dissipated in the pull-up network, independent
of its resistance. The stored $\tfrac12\CL\VDD^2$ is dissipated in the
pull-down network on discharge \cite{rabaey2003}. With switching activity
$\alpha_{\mathrm{sw}}$ and frequency $f$,
\begin{equation}
  E_{\mathrm{CMOS}} = \alpha_{\mathrm{sw}}\,\CL\VDD^2, \qquad
  P_{\mathrm{dyn}} = \alpha_{\mathrm{sw}}\,\CL\VDD^2 f .
  \label{eq:cmos}
\end{equation}

\subsection{Adiabatic Charging}
\label{sec:adiabatic}
If the same load is charged through an on-resistance $R$ by a supply that ramps
linearly from $0$ to $\VDD$ in time $T$, the charging current is approximately
$\CL\VDD/T$, and for $T \gg R\CL$ the dissipated energy is
\cite{athas1994}
\begin{equation}
  E_{\mathrm{ad}} \approx \frac{R\CL}{T}\,\CL\VDD^2 .
  \label{eq:adiabatic}
\end{equation}
This is per transition; a full charge--discharge cycle costs twice this. The
loss falls as $1/T$, and the stored energy is returned to the supply when the
ramp falls. $R$ is the effective on-resistance of the charging path, $T$ the
ramp (transition) time, and $\CL$ the switched capacitance.

For quasi-adiabatic families such as ECRL, the recovery transistor turns off
once the voltage across it drops below its threshold, leaving a residue that
does not scale with $T$ \cite{moon1996}:
\begin{equation}
  E_{\mathrm{ECRL}} \approx \frac{R\CL}{T}\,\CL\VDD^2
  + \tfrac12\,\CL\,|V_{\mathrm{th},p}|^2 .
  \label{eq:ecrl}
\end{equation}
In practice $|V_{\mathrm{th},p}|$ is replaced by the measured residual
voltage, which subthreshold conduction lowers at long $T$ (0.53~V at 10~ns,
0.42~V at 100~ns here), so the residue term itself shrinks slowly with $T$.
\guide{Check the prefactor against \cite{moon1996}.}

\subsection{Energy--Frequency Relationship}
Leakage adds a term that grows with the time a node is held, so the total
energy of a fully adiabatic stage has the form
\begin{equation}
  E(T) = \frac{a}{T} + bT, \qquad
  T^{*} = \sqrt{a/b}, \qquad E_{\min} = 2\sqrt{ab},
  \label{eq:ucurve}
\end{equation}
where $a$ collects the resistive term of \cref{eq:adiabatic} and $b$ the
leakage power. For a symmetric trapezoidal power clock with four equal
intervals of duration $\Tph$, the operating frequency is $f = 1/(4\Tph)$.
Two frequencies follow from the same clock: the transition frequency
$f_{\mathrm{tr}}=1/\Tph$ and the operation rate $f_{\mathrm{clk}}=1/(4\Tph)$.
Plots here use $f_{\mathrm{tr}}$ unless stated. CMOS energy per operation is,
to first order, independent of $f$ until leakage dominates at low $f$.

\subsection{Performance Metrics}
\begin{align}
  \Eop  &= \frac{1}{N}\int_{t_0}^{t_0+NT_{\mathrm{clk}}}
           \sum_k V_k(t)\,I_k(t)\,dt, \label{eq:eop}\\
  P_{\mathrm{avg}} &= \Eop / T_{\mathrm{clk}}, \label{eq:pavg}\\
  \EDP  &= \Eop \times t_{d}, \label{eq:edp}
\end{align}
where the sum runs over every supply or power-clock source $k$, $I_k$ is the
signed current delivered by source $k$ (so recovered energy is subtracted),
$N$ is the number of averaged cycles and $t_d$ the propagation delay.

% =====================================================================
\section{Simulation Methodology}
\label{sec:method}

\subsection{Simulation Environment}
\begin{table}[!t]
\caption{Simulation Environment}
\label{tab:env}
\centering
\begin{tabular}{@{}lL{0.62\columnwidth}@{}}
\toprule
\textbf{Parameter} & \textbf{Value} \\
\midrule
Simulator            & ngspice-42, batch mode, SPARSE 1.3 solver \\
Technology           & SkyWater SKY130 \cite{skywater}, \texttt{tt} corner \\
MOS models           & \texttt{sky130\_fd\_pr\_\_nfet\_01v8}, \texttt{sky130\_fd\_pr\_\_pfet\_01v8} (reduced include, verified identical to the full \texttt{.lib tt}) \\
Temperature          & 27~\si{\celsius} \\
$\VDD$ / clock swing & 1.8~\si{\volt}, 0 to 1.8~V trapezoid \\
Time step            & ECRL: 10~ps; 2LAL: $\Tph/500$ \\
Simulation time      & ECRL: $20\Tph$ (5 clock periods); 2LAL: $8T_{\mathrm{rail}}$ \\
Solver options       & ECRL: defaults; 2LAL: \texttt{reltol=1e-4}, \texttt{abstol=1e-13}, \texttt{vntol=1e-7}, \texttt{chgtol=1e-16}, \texttt{gmin=1e-14}, Gear order 2, \texttt{uic} with rest-state initial conditions \\
\bottomrule
\end{tabular}
\end{table}
\Cref{tab:env} lists the simulation setup. Devices are written in
microns because the models set \texttt{.option scale=1.0u}. The reduced
model include was checked against the full library at four grid points
(both give $1.80491\times10^{-14}$~J for the original ECRL deck).

\subsection{Circuits Under Investigation}
Three circuits are intended: a static CMOS inverter (\cref{fig:cmos}), an
ECRL inverter (\cref{fig:ecrl}) and a 2LAL quad-rail inverter
(\cref{fig:2lal}). \tbd[: the CMOS deck is not yet in the repository].
ECRL is simulated as a single gate driven by ideal slow-ramp inputs. 2LAL is
simulated as a chain: stage 0 drives, stages 1--4 are metered, and stages
5--7 are unmetered loads. Netlists are hand-written and parameterised; the
2LAL cell follows the 4lc PHAS2 cell \cite{debenedictis4lc}.

\begin{figure}[!t]
\centering
\begin{circuitikz}[scale=0.8, transform shape]
  \draw (0,2) node[pmos](Mp){};
  \draw (0,0) node[nmos](Mn){};
  \draw (Mp.S) node[vcc]{$\VDD$};
  \draw (Mn.S) node[ground]{};
  \draw (Mp.D) -- (Mn.D);
  \draw (Mp.G) -- (Mn.G);
  \draw ($(Mp.G)!0.5!(Mn.G)$) node[circ]{} -- ++(-0.6,0) node[left]{$\mathit{in}$};
  \draw ($(Mp.D)!0.5!(Mn.D)$) node[circ]{} -- ++(1.2,0) coordinate(o)
        node[right]{$\mathit{out}$};
  \draw (o) node[circ]{} to[C, l=$\CL$] ++(0,-1.4) node[ground]{};
\end{circuitikz}
\caption{Reference static CMOS inverter.}
\label{fig:cmos}
\end{figure}

\subsection{Common Simulation Conditions}
\begin{table}[!t]
\caption{Common Simulation Conditions}
\label{tab:cond}
\centering
\begin{tabular}{@{}L{0.27\columnwidth}L{0.3\columnwidth}L{0.33\columnwidth}@{}}
\toprule
\textbf{Parameter} & \textbf{Nominal} & \textbf{Swept range} \\
\midrule
$\VDD$               & 1.8~V         & -- \\
$W_n = W_p$          & 1.0~\textmu m  & 0.84, 1.0, 1.2, 1.5, 2.0, 3.0 \\
$L$                  & 0.15~\textmu m & 0.15, 0.18, 0.25, 0.35, 0.5, 0.75, 1.0, 1.5, 2.0 \\
$\CL$ (each metered node) & 25~fF    & 5, 10, 25, 50, 100, 200 \\
$\Tph$               & 10~ns          & 1, 2, 5, 10, 20, 50, 100 \\
$f_{\mathrm{tr}}=1/\Tph$ & 100~MHz    & 1000 to 10 \\
$f_{\mathrm{clk}}=1/(4\Tph)$ & 25~MHz & 250 to 2.5 \\
Input pattern        & 1,0,1,0 (one change per clock period) & -- \\
Input edge           & ECRL: ramp over one full phase in the wait window; 2LAL: rail-derived pulses & -- \\
Stages               & ECRL: 1; 2LAL: 8 (4 metered) & -- \\
\bottomrule
\end{tabular}
\end{table}
Two conditions are not equal by construction. First, the ECRL input
energy is supplied by an ideal source and is not metered, while a 2LAL column
also pays for driving the next stage's gates. Second, per operation an ECRL
gate swings one 25~fF node, whereas the active 2LAL sub-chain swings two
(its $T$ node up and down, its $C$ node down and up).

\subsection{Energy Measurement Method}
\label{sec:meas}
Energy is the integral of the power delivered by the power clock(s), taken
over whole clock periods after start-up, so that energy returned to the clock
is subtracted (\cref{eq:eop}). ngspice counts source current as positive into
the positive terminal, so delivered power is $-V I$.

\emph{ECRL.} A behavioural source converts $-V(\phi)I(V_\phi)$ into a node
voltage, which is integrated over cycle~3 ($[12,16]\Tph$, data $0\to1$,
$E_A$) and cycle~4 ($[16,20]\Tph$, data $1\to0$, $E_B$); $E_C=(E_A+E_B)/2$.
The first three cycles are discarded, and the same two measurements two
periods earlier serve as a settling check:
\begin{lstlisting}
BPCPOWER PCLK 0 V={-V(PHI)*I(VPHI)}
.meas tran E_A INTEG V(PCLK) FROM={12*TPHASE} TO={16*TPHASE}
.meas tran E_B INTEG V(PCLK) FROM={16*TPHASE} TO={20*TPHASE}
\end{lstlisting}

\emph{2LAL.} A 0~V ammeter sits on each rail feeding stages 1--4. A
behavioural current source equal to $\sum_k V(r_k)I(V_{mk})$ charges a 1~F
capacitor, so its voltage is the running energy in joules. Rail periods 4
and 5 are metered and divided by four columns per period. Moving the window
to periods 6--7 changes the result by less than 0.01\%.

\emph{Why the window matters.} The first ECRL deck metered a cycle with no
data transition in it. Holding the data static in the corrected deck
reproduces that value exactly (5.36~fJ at $\Tph=100$~ns); with alternating
data the same point is 12.46~fJ. In both families $E_A$ and $E_B$ agree to
the last printed digit, so no activity correction is applied.

\subsection{Chain Termination (2LAL)}
\label{sec:termination}
The committed 2LAL deck ended at stage~5, whose outputs (node~6) were neither
loaded nor restored. Node~6 then drifts to states in which both sub-chains
read TRUE. This corrupts node~5 through stage~5's reverse pair, and node~4
through stage~4's reverse pair, which is metered. At $\Tph=10$~ns the FALSE
pair at node~4 sits at (0.42, 1.32)~V instead of (0, 1.8)~V. At 50~ns it
regenerates into a spurious TRUE, and the energy doubles (23.97 to 45.11~fJ
between 20 and 50~ns at $W=1$~\textmu m). The original functional checks
sampled $n_1, n_2$ only, and $a_4$ already showed
$(0.13, 1.64)$~V at 10~ns. Adding two unmetered, loaded stages (6 and 7)
removes the defect: node~4 reads $(-0.01, 1.80)$~V at every point, and the
energy becomes monotonic in $\Tph$. A third extra stage changes the result by
0.2\% at 10~ns and 9\% at 100~ns (0.85 vs.\ 0.93~fJ), which bounds the
residual boundary sensitivity. All 2LAL results below use the terminated deck.

\subsection{Functional Verification}
\begin{figure}[!t]
\centering
\simfig{\columnwidth}{figs/fig03_cmos_wave.pdf}\\[2pt]
\simfig{\columnwidth}{figs/fig_ecrl_waveform.png}\\[2pt]
\simfig{\columnwidth}{figs/fig05_2lal_wave.pdf}
\caption{Transient waveforms: (top) CMOS, (middle) ECRL, (bottom) 2LAL.}
\label{fig:waves}
\end{figure}
Functional correctness is checked with \texttt{.meas} samples mid-hold at
both bit values, at every $\Tph$, before any energy is accepted. For ECRL at
$W=1$~\textmu m and $\Tph=10$~ns, $\mathit{in}=1$ gives
$(\mathit{out},\outb)=(0.2~\text{\textmu V}, 1.80~\text{V})$ and $\mathit{in}=0$
the reverse. For 2LAL, TRUE $=(1.8, 0)$ and FALSE $=(0, 1.8)$, and
stages alternate ($a_1=\overline{b}$, $a_2=b$, \dots). In the terminated
deck all sampled nodes, including $n_3$ and $n_4$, are within 20~mV of
the ideal levels (32 checks per run; largest error 23~mV, at $W=3$~\textmu m,
$\Tph=1$~ns), checked at $W=0.84$, 1.0 and 3.0~\textmu m across the $\Tph$ range.
\tbd[: CMOS waveform; export ECRL/2LAL plots to \texttt{figs/}.]

\subsection{Parameter Sweep Methodology}
Five one-at-a-time sweeps are performed about the nominal point of
\cref{tab:cond}: frequency, $\Tph$, transistor width, channel length and load
capacitance. The length and load sweeps are each crossed with the full
$\Tph$ grid.
A two-dimensional sweep of $f$ and $\CL$ produces the operating-region map of
\cref{sec:region} (\tbd, not yet run). Sweeps are scripted
(\texttt{run\_2lal\_sweep.sh}, \texttt{run\_ecrl\_slowramp\_sweep.sh}): each
point is a separate batch run, results pass from \texttt{.meas} to CSV to
gnuplot without manual editing, and a run with no $E_C$ line is treated as
an error. In the length and load sweeps
(\texttt{run\_2lal\_lc\_sweep.sh}, \texttt{run\_ecrl\_lc\_sweep.sh}), a point
whose logic fails (a long channel or a large load that cannot reach the rails
at the fastest ramps) is recorded as a result and excluded from fits and
plots. A failed substitution, a simulator error, or a $W$, $L$ or $\Tph$
read back from the simulator that differs from the request still stops the
sweep.

\subsection{Choice of the Fixed Operating Point}
\label{sec:fixedpoint}
The length sweep holds $W=1.0$~\textmu m and $\CL=25$~fF; the load sweep holds
$W=1.0$~\textmu m and $L=0.15$~\textmu m. The reasons are:
\begin{itemize}
\item \emph{It is the anchor of every other result.} Both sweeps pass
  through $W=1.0$~\textmu m, $L=0.15$~\textmu m, $\CL=25$~fF, and each sweep
  script checks at the end that it reproduces the width sweep there; it does,
  to every printed digit. A change in a deck or in the models therefore shows
  up as a failed anchor.
\item \emph{A formula-optimal width does not stay optimal across a length
  sweep.} The adiabatic term behaves as $a\propto R\,(\CL+kW)^2$ with
  $R\propto1/W$, which is minimised at $W^{*}=\CL/k$: widening pays only while
  the transistor's own capacitance stays below the load \cite{fischer2003,
  fischer2004}. At $L=0.15$~\textmu m the width sweep gives
  $k=7.1$~fF/\textmu m, so $W^{*}=3.5$~\textmu m. The gate part of $k$ grows
  with $L$ (12.6~fF/\textmu m per \textmu m of $L$, \cref{sec:length}), so
  $k=7.1+12.6(L-0.15)$ and $W^{*}$ falls to about 1.4~\textmu m at
  $L=1$~\textmu m and 0.8~\textmu m at $L=2$~\textmu m. Fixing $W$ at
  3.5~\textmu m would turn the long-channel points into a measurement of gate
  capacitance.
\item \emph{The extremes are worse.} Minimum width (0.42~\textmu m)
  maximises $R$, so long channels fail logic at even more ramp times, and it
  sits beside narrow-width model-bin edges. The largest width lets the gate
  load, which grows as $W L$, dominate $\CL$.
\item \emph{$L=0.15$~\textmu m for the load sweep} is the minimum length,
  which the length sweep shows to be the lowest-energy length at every $\Tph$
  except 100~ns, where 0.18~\textmu m is 0.8\% lower because it leaks less.
\end{itemize}
The 2LAL deck previously set the source/drain diffusion area as
$W\cdot L\cdot1.77$, so a longer channel also grew the junctions. Diffusion
geometry now uses a separate $L_{\mathrm{dif}}=0.15$~\textmu m. At
$L=0.15$~\textmu m the deck is numerically unchanged.

% =====================================================================
\section{Experimental Investigation}
\label{sec:results}
All energies are $E_C$ in fJ per operation (one operation = one clock period
per ECRL gate, or per quad-rail 2LAL column). ECRL values are from
\texttt{ecrl\_slowramp\_sweep.csv}; 2LAL values are from the terminated deck
(\cref{sec:termination}). No CMOS deck has been simulated yet; where a CMOS
reference is needed, the analytical dissipation per transition,
$\tfrac12\CL\VDD^2=40.5$~fJ for 25~fF, is used and labelled as such.

\subsection{Baseline Comparison}
Fixed conditions: $\VDD=1.8$~V, $W=1.0$~\textmu m, $L=0.15$~\textmu m,
$\CL=25$~fF, $\Tph=10$~ns ($f_{\mathrm{tr}}=100$~MHz, $f_{\mathrm{clk}}=25$~MHz).
\begin{table}[!t]
\caption{Baseline Comparison ($W=1$~\textmu m, $\Tph=10$~ns)}
\label{tab:baseline}
\centering
\setlength{\tabcolsep}{4pt}
\begin{tabular}{@{}lcccc@{}}
\toprule
\textbf{Logic} & \textbf{$E$ (fJ)} & \textbf{$P$ (\si{\micro\watt})}
  & \textbf{$t_d$ (ns)} & \textbf{EDP (fJ$\cdot$ns)} \\
\midrule
CMOS & \tbd\,$^{a}$ & \tbd & \tbd & \tbd \\
ECRL & 21.66 & 0.542 & 10 & 216.6 \\
2LAL & 6.31  & 0.158 & 10 & 63.1 \\
\bottomrule
\multicolumn{5}{@{}p{0.95\columnwidth}@{}}{\footnotesize
$P=E\,f_{\mathrm{clk}}$. $t_d$ is the clock-defined latency of one stage
($\Tph$); see \cref{sec:edp}. $^{a}$Analytical reference
$\tfrac12\CL\VDD^2=40.5$~fJ per transition.}
\end{tabular}
\end{table}
\begin{figure}[!t]
\centering
\simfig{\columnwidth}{figs/fig06_energy_baseline.pdf}
\caption{Energy per operation at the nominal point.}
\label{fig:baseline}
\end{figure}
At the nominal point, 2LAL dissipates 6.31~fJ and ECRL 21.66~fJ per
operation, a ratio of 3.43. Both are below the 40.5~fJ analytical CMOS
reference, by 84\% (2LAL) and 47\% (ECRL). The ECRL gate returns 62.9\% of
the energy its clock delivers (36.72 of 58.38~fJ), and the 2LAL column returns
90.8\% (62.28 of 68.58~fJ).

\subsection{Effect of Frequency}
\label{sec:freq}
$\Tph$ is swept from 1 to 100~ns with all other parameters at nominal.
\begin{table}[!t]
\caption{Energy per Operation (fJ) Versus Frequency ($W=1$~\textmu m)}
\label{tab:freq}
\centering
\setlength{\tabcolsep}{3.5pt}
\begin{tabular}{@{}cccccc@{}}
\toprule
\textbf{$f_{\mathrm{tr}}$} & \textbf{$f_{\mathrm{clk}}$} & \textbf{CMOS}
  & \textbf{ECRL} & \textbf{2LAL} & \textbf{Ratio} \\
(MHz) & (MHz) & (fJ) & (fJ) & (fJ) & E/2 \\
\midrule
1000 & 250  & \tbd & 54.69 & 46.13 & 1.19 \\
500  & 125  & \tbd & 39.82 & 26.45 & 1.51 \\
200  & 50   & \tbd & 27.32 & 11.93 & 2.29 \\
100  & 25   & \tbd & 21.66 & 6.31  & 3.43 \\
50   & 12.5 & \tbd & 17.83 & 3.31  & 5.38 \\
20   & 5    & \tbd & 14.37 & 1.39  & 10.3 \\
10   & 2.5  & \tbd & 12.46 & 0.85  & 14.6 \\
\bottomrule
\multicolumn{6}{@{}l@{}}{\footnotesize Ratio = ECRL/2LAL.} \\
\end{tabular}
\end{table}
\begin{figure}[!t]
\centering
\simfig{\columnwidth}{figs/fig07_energy_vs_freq.pdf}\\[2pt]
\simfig{\columnwidth}{figs/fig08_power_vs_freq.pdf}
\caption{(top) Energy and (bottom) average power versus $f_{\mathrm{tr}}$.}
\label{fig:freq}
\end{figure}
\emph{(1) How energy changes with frequency.} Energy per operation rises
monotonically with frequency in both families. From 10 to 1000~MHz, ECRL
rises 4.4$\times$ (12.46 to 54.69~fJ) and 2LAL rises 54$\times$ (0.85 to
46.13~fJ). Average power $E f_{\mathrm{clk}}$ rises much faster because both
factors grow: ECRL goes from 31~nW to 13.7~\textmu W, and 2LAL from 2.1~nW to
11.5~\textmu W.

\emph{(2) Strongest frequency dependence.} 2LAL. Its log--log slope of $E$
against $\Tph$ is $-0.94$ between 10 and 50~ns at $W=1$~\textmu m ($-0.89$ to
$-0.99$ over all widths). ECRL's slope is only $-0.24$ from 10 to 100~MHz and
$-0.40$ from 100 to 1000~MHz.

\emph{(3) Where the relative advantage changes.} With the terminated deck,
the families do not cross within 10--1000~MHz: 2LAL is lower at every point.
The gap narrows towards high frequency (1.19$\times$ at 1000~MHz). Against the
analytical CMOS reference, ECRL reaches 40.5~fJ between 500~MHz (39.82~fJ)
and 1000~MHz (54.69~fJ), and 2LAL between 500~MHz (26.45~fJ) and 1000~MHz
(46.13~fJ). \guide{Replace with the simulated CMOS crossover.}

\emph{(4) Agreement with theory.} 2LAL follows the $1/T$ law of
\cref{eq:adiabatic} over 5--50~ns. Its $E\cdot\Tph$ product rises only from
59.7 to 69.5~fJ$\cdot$ns while $\Tph$ grows 10$\times$. The deviation at
1~ns ($E\Tph=46.1$) is expected as $\Tph$ approaches $RC$. ECRL does not
follow a pure $c + a/T$ law: a least-squares fit leaves systematic residuals
of up to 3.7~fJ. An empirical $c + a/\sqrt{T}$ form fits within 0.9~fJ, but
we attach no physical meaning to the exponent. The shallow slope agrees with
\cref{eq:ecrl}, which has a residue term that does not scale as $1/T$. That
residue itself falls slowly with $T$, because the stranded voltage decays
from 0.53~V to 0.42~V between 10 and 100~ns.

\subsection{Effect of Transition Time}
\label{sec:tph}
Because $f_{\mathrm{clk}}=1/(4\Tph)$, this sweep is the same data as
\cref{tab:freq} viewed on the time axis.
\begin{table}[!t]
\caption{Energy (fJ) and Returned Fraction Versus $\Tph$ ($W=1$~\textmu m)}
\label{tab:tph}
\centering
\setlength{\tabcolsep}{4pt}
\begin{tabular}{@{}cccccc@{}}
\toprule
\textbf{$\Tph$} & \textbf{$f_{\mathrm{clk}}$} & \textbf{ECRL} & \textbf{2LAL}
  & \textbf{ECRL} & \textbf{2LAL} \\
(ns) & (MHz) & $E$ (fJ) & $E$ (fJ) & $E\Tph$ & $E\Tph$ \\
\midrule
1   & 250  & 54.69 & 46.13 & 54.7   & 46.1 \\
2   & 125  & 39.82 & 26.45 & 79.6   & 52.9 \\
5   & 50   & 27.32 & 11.93 & 136.6  & 59.7 \\
10  & 25   & 21.66 & 6.31  & 216.6  & 63.1 \\
20  & 12.5 & 17.83 & 3.31  & 356.5  & 66.3 \\
50  & 5    & 14.37 & 1.39  & 718.3  & 69.5 \\
100 & 2.5  & 12.46 & 0.85  & 1245.9 & 85.3 \\
\midrule
\multicolumn{6}{@{}p{0.95\columnwidth}@{}}{\footnotesize Returned/delivered
energy: ECRL 62.9\% (10~ns), 76.0\% (100~ns); 2LAL 90.8\% (10~ns), 98.8\%
(100~ns). $E\Tph$ in fJ$\cdot$ns.} \\
\bottomrule
\end{tabular}
\end{table}
\begin{figure}[!t]
\centering
\simfig{\columnwidth}{figs/fig09_energy_vs_tph.pdf}
\caption{Energy versus $\Tph$ (log--log).}
\label{fig:tph}
\end{figure}
Transition time sets dissipation because the charging current is
$\approx\CL\VDD/\Tph$, and the resistive loss is $I^2R$ integrated over
$\Tph$. Doubling $\Tph$ halves the current, quarters the power and doubles the
duration, so the loss halves (\cref{eq:adiabatic}). 2LAL shows this directly
and returns 98.8\% of the delivered energy at 100~ns. ECRL recovers less
(76.0\% at 100~ns) because part of each cycle's loss is the threshold residue
dumped to ground when the data changes.

The leakage term $b\Tph$ of \cref{eq:ucurve} does not produce a minimum within
the swept range: 2LAL energy still falls from 50 to 100~ns at every width.
It does flatten, however: $E\Tph$ rises from 69.5 to 85.3~fJ$\cdot$ns at
$W=1$~\textmu m, and at $W=0.84$~\textmu m the slope from 50 to 100~ns is
only $-0.5$. Locating $T^{*}$ requires $\Tph>100$~ns. \tbd

\subsection{Effect of Transistor Width}
\begin{table}[!t]
\caption{Energy (fJ) Versus Transistor Width at Two Ramp Times}
\label{tab:width}
\centering
\setlength{\tabcolsep}{4pt}
\begin{tabular}{@{}ccccc@{}}
\toprule
& \multicolumn{2}{c}{\textbf{ECRL}} & \multicolumn{2}{c}{\textbf{2LAL}} \\
\cmidrule(lr){2-3}\cmidrule(l){4-5}
\textbf{$W$ (\textmu m)} & \textbf{1~ns} & \textbf{100~ns} & \textbf{1~ns} & \textbf{100~ns} \\
\midrule
0.84 & 59.59 & 12.61 & 50.53 & 1.16 \\
1.0  & 54.69 & 12.46 & 46.13 & 0.85 \\
1.2  & 51.12 & 12.56 & 42.94 & 0.62 \\
1.5  & 46.05 & 11.97 & 38.27 & 0.55 \\
2.0  & 41.40 & 11.73 & 35.76 & 0.65 \\
3.0  & 36.43 & 11.55 & 33.79 & 0.77 \\
\bottomrule
\end{tabular}
\end{table}
\begin{figure}[!t]
\centering
\simfig{\columnwidth}{figs/fig10_energy_vs_width.pdf}\\[2pt]
\simfig{\columnwidth}{figs/fig11_delay_vs_width.pdf}
\caption{(top) Energy versus width; (bottom) \tbd[: output-to-clock lag
versus width, see text].}
\label{fig:width}
\end{figure}
Widening the devices lowers the channel resistance $R$ in
\cref{eq:adiabatic}, but it adds gate and junction capacitance. The benefit
therefore appears only where the loss is resistive. Going from 0.84 to
3.0~\textmu m lowers energy at 1~ns by 38.9\% (ECRL) and 33.1\% (2LAL), and
at 10~ns by 19.8\% (ECRL, 22.55 to 18.09~fJ) and 39.8\% (2LAL, 6.89 to
4.15~fJ). At 100~ns ECRL improves by only 8.4\%, because its remaining loss is
mostly the threshold residue, which does not depend on $R$. 2LAL at 100~ns has
its minimum at 1.5~\textmu m (0.55~fJ) and rises 40\% by 3.0~\textmu m. This
is consistent with extra capacitance and leakage outweighing the lower
resistance once the $RC/T$ term is small.

In these cells, delay is set by the clock, not by drive strength: each stage
advances one phase ($\Tph$) as long as its output settles within the ramp. A
width sweep therefore changes how closely the output tracks the rail, not the
latency. For example, at $\Tph=10$~ns the ECRL output lags the clock at
mid-ramp (0.24~V against 0.90~V), but at 100~ns it tracks it (0.89~V against
0.90~V). \tbd[: add a lag or settling metric versus $W$ if the delay figure is kept.]
\guide{The committed 2LAL sweep showed wider devices as worse at every
frequency. That trend was produced by the termination defect and should not
be reported.}

\subsection{Effect of Load Capacitance}
\label{sec:cl}
Theory predicts that CMOS scales as $\CL$ (\cref{eq:cmos}), the adiabatic
resistive term as $\CL^2$ at fixed $R$ and $T$ (\cref{eq:adiabatic}), and the
ECRL residue as $\CL$ (\cref{eq:ecrl}). $\CL$ is swept from 5 to 200~fF at
$W=1.0$~\textmu m, $L=0.15$~\textmu m (\cref{tab:cl,fig:cl}).
\begin{table}[!t]
\caption{Energy (fJ) Versus Load Capacitance ($W=1$~\textmu m, $L=0.15$~\textmu m)}
\label{tab:cl}
\centering
\setlength{\tabcolsep}{4pt}
\begin{tabular}{@{}ccccc@{}}
\toprule
& \multicolumn{2}{c}{\textbf{ECRL}} & \multicolumn{2}{c}{\textbf{2LAL}} \\
\cmidrule(lr){2-3}\cmidrule(l){4-5}
\textbf{$\CL$ (fF)} & \textbf{10~ns} & \textbf{100~ns} & \textbf{10~ns} & \textbf{100~ns} \\
\midrule
5   & 3.12   & 2.00   & 0.95   & 0.38 \\
10  & 6.96   & 4.30   & 1.85   & 0.47 \\
25  & 21.66  & 12.46  & 6.29   & 0.93 \\
50  & 53.90  & 28.48  & 19.10  & 2.34 \\
100 & 140.54 & 66.44  & 62.98  & 7.49 \\
200 & 381.24 & 159.15 & 212.96 & 26.86 \\
\bottomrule
\end{tabular}
\end{table}
\begin{figure}[!t]
\centering
\simfig{\columnwidth}{../plots/2lal_energy_vs_cap.png}\\[2pt]
\simfig{\columnwidth}{../plots/ecrl_energy_vs_cap.png}
\caption{Energy versus load capacitance, one curve per $\Tph$: (top) 2LAL,
(bottom) ECRL. Dashed lines are $E\propto\CL^2$ and $E\propto\CL$ guides.
Points whose logic fails are omitted.}
\label{fig:cl}
\end{figure}
The two families separate as predicted. For 2LAL, fitting
$E(\Tph)$ at each load and writing the $1/\Tph$ coefficient as
$a=4R(\CL+C_p)^2\VDD^2$ (four node swings per operation) makes
$\sqrt{a/4\VDD^2}$ a straight line in $\CL$; it is linear within 1.0\% at
all six loads. Its slope gives $R=4.54$~k$\Omega$, against 4.92~k$\Omega$
measured directly as the mean DC on-resistance of the transmission gate
($-8\%$), and its intercept gives $C_p=7.7$~fF. Between 25 and 200~fF each
doubling of $\CL$ multiplies 2LAL energy at 10~ns by 3.0--3.4 (log--log slope
1.6--1.8), approaching the $\times4$ of a pure $\CL^2$ law. For ECRL the fixed
residue grows from 2.5 to 153~fJ across the sweep, and at 100~ns energy is
close to linear in $\CL$ (log--log slope 1.2--1.3): the threshold loss of
\cref{eq:ecrl} dominates. ECRL and 2LAL fail logic at $\Tph=1$~ns from
100~fF, and at 2~ns at 200~fF, where $R\CL$ approaches the ramp time.

\subsection{Effect of Channel Length}
\label{sec:length}
$L$ is swept from 0.15 to 2.0~\textmu m at $W=1.0$~\textmu m, $\CL=25$~fF
(\cref{tab:length,fig:length}). The grid includes SKY130's characterised
lengths (0.15, 0.18, 0.25, 0.5, 1 and 2~\textmu m), where model bins meet.
\begin{table}[!t]
\caption{Energy (fJ) Versus Channel Length ($W=1$~\textmu m, $\CL=25$~fF)}
\label{tab:length}
\centering
\setlength{\tabcolsep}{4pt}
\begin{tabular}{@{}ccccc@{}}
\toprule
& \multicolumn{2}{c}{\textbf{ECRL}} & \multicolumn{2}{c}{\textbf{2LAL}} \\
\cmidrule(lr){2-3}\cmidrule(l){4-5}
\textbf{$L$ (\textmu m)} & \textbf{10~ns} & \textbf{100~ns} & \textbf{10~ns} & \textbf{100~ns} \\
\midrule
0.15 & 21.66 & 12.46 & 6.29   & 0.93 \\
0.18 & 25.43 & 15.00 & 8.48   & 0.92 \\
0.25 & 31.31 & 18.75 & 13.49  & 1.63 \\
0.5  & 41.75 & 24.17 & 28.01  & 4.17 \\
1.0  & 53.17 & 28.31 & 59.21  & 10.93 \\
2.0  & 67.20 & 33.88 & fails  & 34.87 \\
\bottomrule
\end{tabular}
\end{table}
\begin{figure}[!t]
\centering
\simfig{\columnwidth}{../plots/2lal_energy_vs_length.png}\\[2pt]
\simfig{\columnwidth}{../plots/ecrl_vs_2lal_length_T10.png}
\caption{(top) 2LAL energy versus channel length, one curve per $\Tph$;
(bottom) ECRL and 2LAL at $\Tph=10$~ns.}
\label{fig:length}
\end{figure}
Energy rises with $L$ in both families, and the shortest channel is the
cheapest at every $\Tph$ except 100~ns, where 2LAL at 0.18~\textmu m is
0.8\% below 0.15~\textmu m because the longer channel leaks less. 2LAL rises
faster: at 10~ns it is 3.4$\times$ below ECRL at 0.15~\textmu m but costs
more above $L\approx0.85$~\textmu m. Its $1/\Tph$ coefficient grows as
$R(L)\,C_{\mathrm{eff}}(L)^2$, and both factors grow with $L$. With the
transmission-gate resistance $R(L)$ measured at each length (4.9 to
86~k$\Omega$, proportional to $L$ above 0.5~\textmu m), the effective node
capacitance $C_{\mathrm{eff}}=\sqrt{a/4\VDD^2R}$ is
$28.2+12.6\,L$~fF (\textmu m), linear within 5.3\%: the 25~fF load plus
junctions, plus gate charge growing as $WL$. The slope lies between two
oxide-based estimates ($C_{ox}=8.3$~fF/\textmu m$^2$ from the model's
$t_{ox}=4.15$~nm): 16.6~fF/\textmu m for the two gates per node that are in
inversion while the node swings, and about 10~fF/\textmu m if those gates
are counted only above threshold. For ECRL the threshold residue grows from
15.1 to 27.6~fJ between 0.15 and 1~\textmu m, following the pull-up pFET,
whose $|V_{\mathrm{th}}|$ rises from 0.73 to 0.94~V. The nFET threshold falls
over the same range (0.68 to 0.55~V), the reverse short-channel effect of
SKY130's halo implants. Long channels fail logic at the fastest ramps
(2LAL from 0.35~\textmu m at 1~ns; ECRL from 0.35~\textmu m at 1~ns).

\subsection{Energy--Delay Analysis}
\label{sec:edp}
With $t_d=\Tph$ per stage, $\EDP=E\,\Tph$ (\cref{tab:tph}, last two
columns). EDP rises with $\Tph$ in both families, so within this sweep EDP is
minimised at the shortest ramp: 46.1~fJ$\cdot$ns for 2LAL and
54.7~fJ$\cdot$ns for ECRL at 1~ns. The penalty for slowing down differs
greatly. From 1 to 100~ns, ECRL's EDP grows 22.8$\times$ while its energy
falls only 4.4$\times$; 2LAL's EDP grows 1.85$\times$ while its energy falls
54$\times$. In other words, slowing 2LAL buys energy almost one-for-one in
time, whereas slowing ECRL mostly costs time.
\begin{table}[!t]
\caption{Energy-Delay Product (fJ$\cdot$ns), $W=1$~\textmu m, $t_d=\Tph$}
\label{tab:edp}
\centering
\begin{tabular}{@{}lccc@{}}
\toprule
\textbf{Condition} & \textbf{CMOS} & \textbf{ECRL} & \textbf{2LAL} \\
\midrule
$\Tph=1$~ns   & \tbd & 54.7   & 46.1 \\
$\Tph=10$~ns  & \tbd & 216.6  & 63.1 \\
$\Tph=100$~ns & \tbd & 1245.9 & 85.3 \\
\bottomrule
\end{tabular}
\end{table}
\begin{figure}[!t]
\centering
\simfig{\columnwidth}{figs/fig13_edp.pdf}
\caption{Energy-delay product comparison.}
\label{fig:edp}
\end{figure}

\subsection{Comprehensive Comparison}
\begin{table}[!t]
\caption{Comprehensive Comparison ($W=1$~\textmu m, $\Tph=10$~ns)}
\label{tab:summary}
\centering
\setlength{\tabcolsep}{3pt}
\begin{tabular}{@{}L{0.30\columnwidth}L{0.14\columnwidth}L{0.22\columnwidth}L{0.24\columnwidth}@{}}
\toprule
\textbf{Parameter} & \textbf{CMOS} & \textbf{ECRL} & \textbf{2LAL} \\
\midrule
Energy/op.\ (fJ)          & \tbd & 21.66 & 6.31 \\
Avg.\ power (\si{\micro\watt}) & \tbd & 0.542 & 0.158 \\
Latency/stage (ns)        & \tbd & 10 & 10 \\
EDP (fJ$\cdot$ns)         & \tbd & 216.6 & 63.1 \\
Transistors               & 2 & 4 per gate & 8 per cell; 16 per inverter column \\
Clock rails               & DC $\VDD$ & 1 per gate; 4 phases in a cascade & 4 \\
Encoding                  & single-rail & dual-rail & dual-rail; quad-rail to invert \\
Nodes swung per op.       & 1 & 1 & 2 \\
Returned / delivered      & 0 & 62.9\% & 90.8\% \\
Non-adiabatic residue     & -- & yes (\cref{eq:ecrl}) & none observed \\
\bottomrule
\end{tabular}
\end{table}
\Cref{tab:summary} collects the nominal-point comparison. The lower energy of
2LAL comes with 4$\times$ the transistors per inverter, four clock rails, and
twice the swung capacitance per operation. Its standing capacitance, area and
leakage are all larger.

\subsection{Operating-Region Analysis}
\label{sec:region}
\tbd[: requires the $f\times\CL$ sweep and a simulated CMOS reference.] Along
the frequency axis at $\CL=25$~fF, the ratio to the analytical CMOS reference
is $\rho_{\mathrm{ECRL}}=0.31$ at 10~MHz and $1.35$ at 1000~MHz, and
$\rho_{\mathrm{2LAL}}=0.021$ at 10~MHz and $1.14$ at 1000~MHz. For both
families, $\rho<1$ holds up to at least 500~MHz and is lost by 1000~MHz.
\begin{figure*}[!t]
\centering
\simfig{0.48\textwidth}{figs/fig14a_region_ecrl.pdf}\hfill
\simfig{0.48\textwidth}{figs/fig14b_region_2lal.pdf}
\caption{Energy ratio $\rho$ versus frequency (x) and load capacitance (y) for
(left) ECRL and (right) 2LAL; the $\rho=1$ contour is marked.}
\label{fig:region}
\end{figure*}

\section{Discussion}
\label{sec:discussion}
\emph{Why adiabatic energy depends on transition time.} The loss is $I^2R$
over the ramp, and $I\propto1/\Tph$, so the loss scales as $R\CL/\Tph$
(\cref{eq:adiabatic}). 2LAL shows this slope ($-0.89$ to $-0.99$). ECRL
shows a much weaker one because a part of its loss does not scale with
$\Tph$.

\emph{Why frequency influences dissipation.} Frequency and ramp time are the
same knob ($f_{\mathrm{clk}}=1/4\Tph$). Per operation, energy rises with $f$.
Per second, power rises faster still, because the operation rate rises too.

\emph{Why load capacitance matters.} $\CL$ sets both the charge moved per
transition and the $RC$ time constant, so the resistive term scales as
$\CL^2$ at fixed $R$ and $T$. The load sweep confirms this for 2LAL, with
an $R$ and a parasitic $C_p$ that match independent measurements, and shows
the ECRL threshold residue growing close to linearly with $\CL$
(\cref{sec:cl}).

\emph{Why channel length matters.} A longer channel raises the charging-path
resistance roughly in proportion to $L$ \cite{fischer2004} and adds gate
capacitance in proportion to $WL$; the adiabatic term grows as the product
$R\,C^2$, so energy per operation rises steeply with $L$ in 2LAL. Longer
channels leak less, but in this range that saving is visible only at the
slowest ramp (\cref{sec:length}).

\emph{Why width affects energy and delay.} Width trades lower $R$ against
higher device capacitance and leakage. It helps where the loss is resistive:
short ramps in both families, and 2LAL generally. It does little where the
loss is a threshold residue (ECRL at long ramps). It hurts where the resistive
term is already small (2LAL above 1.5~\textmu m at 100~ns). Stage latency is
fixed by the clock, so width affects how closely the output tracks the rail,
not the delay.

\emph{Practical differences between ECRL and 2LAL.} ECRL needs 4 FETs and one
clock per gate but strands $\approx|V_{\mathrm{th},p}|$ on the released
output. Its per-cycle energy also depends on data activity: 5.36~fJ with
static data against 12.46~fJ with alternating data at 100~ns. 2LAL needs
8~FETs per cell, 16 per inverter column and four rails. It has no residue, but
its floating FALSE nodes make it sensitive to what surrounds the metered
stages (\cref{sec:termination}). In this study, that sensitivity changed the
headline comparison.

\emph{Comparison with published results.} The ECRL trend reported in
\cite{moon1996}, energy rising with frequency and not extrapolating to zero,
is reproduced qualitatively. The numbers are not comparable because the
process, supply and device sizes differ. For 2LAL, Frank \emph{et al.} report
6.9~aJ per cycle per FET for a shift register at 1~MHz, $\CL=10$~fF,
$W=480$~nm in Sandia's 180~nm process \cite{frank2020talk}. Simulated at the
same point in SKY130 ($\Tph=250$~ns, $W=0.48$~\textmu m, $L=0.18$~\textmu m),
the ring passes all 96 logic checks and dissipates 8.1~aJ per FET per cycle,
1.17$\times$ the published value; a quad-rail column has 16 FETs and one
sub-chain pulsing per cycle, i.e.\ 50\% activity per sub-chain. Functionally,
the ring follows the 2LAL timing rule of one tick of delay per stage
\cite{frank2017tutorial}: the waveform decoder finds every stage's output to
be the inverse of its input one $\Tph$ later (154/154 stage-periods).
\guide{No published ECRL or 2LAL simulation in SKY130 was found; the closest
same-PDK reference is DeBenedictis's ngspice shift-register decks, which
include a SKY130 option \cite{debenedictis4lc}. A one-parameter-at-a-time
sweep of $W$, $\CL$ and frequency is reported in \cite{reddy2023} (full text
not accessed).}

\emph{Conditions that favour energy recovery.} Recovery is favoured by long
ramps relative to $R\CL$, switches that never turn on across a voltage and
never turn off while carrying current (2LAL returns 98.8\% at 100~ns), and,
for ECRL, low data activity, because the residue is lost only when the data
changes.

\subsection{Comparison With Previous Work}
\begin{table*}[!t]
\caption{Comparison With Previously Reported Results}
\label{tab:prior}
\centering
\begin{tabular}{@{}L{0.075\textwidth}L{0.12\textwidth}L{0.09\textwidth}L{0.065\textwidth}L{0.11\textwidth}L{0.09\textwidth}L{0.285\textwidth}@{}}
\toprule
\textbf{Ref.} & \textbf{Technology} & \textbf{Logic} & \textbf{$\VDD$}
  & \textbf{Frequency} & \textbf{Energy} & \textbf{Key finding} \\
\midrule
\cite{moon1996} & 1.0~\textmu m CMOS\,$^{b}$ & ECRL & 5~V\,$^{b}$ & \tbd & \tbd & Energy rises with frequency \\
\cite{frank2020talk} & Sandia MESA 180~nm (sim.) & 2LAL shift reg. & \tbd & 1~MHz & 6.9~aJ/FET/cycle & $\CL=10$~fF, $W=480$~nm, 50\% activity; 37~aJ in MESA 350~nm \\
This work, same point & SKY130 (tt) & 2LAL ring & 1.8~V & 1~MHz & 8.1~aJ/FET/cycle & $\CL=10$~fF, $W=0.48$, $L=0.18$~\textmu m; 1.17$\times$ the published value, all 96 logic checks pass \\
\cite{fischer2004} & 0.13~\textmu m CMOS (sim.) & PFAL & \tbd & 1~MHz--1~GHz & 0.60--0.88~fJ at 500~MHz & Sizing rule: widen only while device capacitance $<$ load \\
This work & SKY130 (tt) & ECRL, 2LAL & 1.8~V & 10--1000~MHz ($f_{\mathrm{tr}}$) & 0.55--59.6~fJ & 2LAL lower at all swept $f$; ECRL flattens to a threshold residue \\
\midrule
\multicolumn{7}{@{}l@{}}{\footnotesize $^{b}$As read from \cite{moon1996} (also $V_t\approx0.2$~V, 5/10~\textmu m devices); verify against the paper.} \\
\bottomrule
\end{tabular}
\end{table*}
\guide{Do not claim no prior study exists; state how this work extends or
systematically evaluates established approaches.}

% =====================================================================
\section{Limitations}
\label{sec:limits}
This is a pre-layout simulation study in one process corner (SKY130 tt,
27~$^\circ$C). It covers inverters only, with no extracted parasitics and no
fabricated or measured silicon. The power-clock generator is ideal, and its
losses are excluded, as in most adiabatic results. No CMOS deck or supply
sweep has been run yet. ECRL is a single gate
driven by ideal slow-ramp inputs whose energy is not metered. The 2LAL cell
is a reconstruction from secondary sources. The lowest 2LAL energies (below
1~fJ) are sensitive to chain termination at the 10\% level.

\section{Conclusion}
\label{sec:conclusion}
This work presented a comparative simulation study of ECRL and 2LAL
inverters in SKY130 under one simulator, one load and one energy definition.
The influence of transition time (equivalently frequency) and transistor
width was investigated at $\VDD=1.8$~V and $\CL=25$~fF. 2LAL dissipated less
than ECRL at every swept point, by 1.19$\times$ at 1~ns and 14.6$\times$ at
100~ns ($W=1$~\textmu m), and it followed the adiabatic $1/T$ law. ECRL
flattened towards a threshold-related residue and depended on data activity.
The lowest observed energy was 0.55~fJ (2LAL, $W=1.5$~\textmu m,
$\Tph=100$~ns), and the lowest ECRL energy was 11.55~fJ ($W=3$~\textmu m,
$\Tph=100$~ns). The EDP analysis showed that slowing 2LAL buys energy almost
one-for-one in time, whereas slowing ECRL mostly costs time. Energy rose
with channel length in both families, steeply enough in 2LAL that it lost its
advantage above $L\approx0.85$~\textmu m at 10~ns, and the load sweep separated
the $\CL^2$ adiabatic term from ECRL's near-linear threshold residue. Overall, the
study shows that adiabatic energy figures depend on the metering window and
on the test-structure boundaries as much as on the circuit: two such defects
each reversed a conclusion here.
\tbd[: CMOS comparison sentence once simulated.]

\section*{Future Work}
\guide{Keep only realistic directions.}
NAND, NOR and multi-gate chains in both families. An ECRL chain in place of
the ideal-input single gate. Separation of leakage from switching energy by
measuring at zero activity. Supply-voltage sweeps.
$\Tph>100$~ns to locate the 2LAL minimum. Layout with extracted parasitics,
where 2LAL's larger standing capacitance will show. Costing the power-clock
generator.

\section*{Acknowledgment}
\tbd.

% =====================================================================
\bibliographystyle{IEEEtran}
\bibliography{refs}

% ---------------- supplementary (optional) ----------------
% \clearpage
% \appendix
% \section{SPICE Decks and Reproduction}
% \lstinputlisting{spice/ecrl_chain.cir}

\end{document}
